% Source: https://artofproblemsolving.com/community/c3h1525482_transformations
% Stable identity: aops-topic-1525482-q01
\textbf{(a)} $f\left(\dfrac {x}{2}\right)=\dfrac{mx}{2}=\dfrac{1}{2}f(x)$ for every $x$, so the two graphs coincide.

\textbf{(b)} Take $f(x)=x+1$. Then $f\left(\dfrac {x}{2}\right)=\dfrac {x}{2}+1$, while $\dfrac{1}{2}f(x)=\dfrac {x}{2}+\dfrac{1}{2}$; at $x=0$ the values are $1$ and $\dfrac{1}{2}$, so the graphs differ. (Another example: for $f(x)=x^2$ the horizontal stretch gives $y=\dfrac{x^2}{4}$, which is a vertical shrink by a factor of $4$, not $2$.) The same examples show that a horizontal shrink is not the same as a vertical stretch either.

\textbf{(c)} Write $f(x)=a_0+a_1x+\cdots+a_nx^n$. The condition says
\[
\sum_{k=0}^{n}a_k\left(\dfrac{1}{2^k}-\dfrac{1}{2}\right)x^k=0\qquad\text{for all }x .
\]
A polynomial that is zero for all real $x$ has all coefficients $0$, so $a_k\left(\dfrac{1}{2^k}-\dfrac{1}{2}\right)=0$ for each $k$. The bracket is nonzero unless $k=1$, so $a_k=0$ for all $k\ne1$, and $f(x)=a_1x$. $\blacksquare$
